% GENERATED FILE. Edit canonical lessons, then run npm run generate:books.
% canonical-source-hash: fnv1a64:a4c8e8b5fc059b71
% canonical-lessons: KA-C222-pyant, KA-C222-langa, KA-C222-ravike, KA-C222-baniyan, KA-C222-butu

\chapter{Trousers, Long skirt, Blouse, Vest, Shoes}
\label{ch:ka-trousers-long-skirt-blouse-vest-shoes}
\hlchaptermodality{222}

\begin{chapteropening}
\textbf{By the end of this chapter:} \emph{I can say trousers, a long skirt, a blouse, a vest and shoes.}

Complete the last lesson of chapter 222: I can say trousers, a long skirt, a blouse, a vest and shoes.
\end{chapteropening}

\section[\texorpdfstring{\kn{ಪ್ಯಾಂಟ್} (pyāṇṭ)}{ಪ್ಯಾಂಟ್ (pyāṇṭ)}]{\kn{ಪ್ಯಾಂಟ್} (pyāṇṭ) — trousers}

\label{lesson:KA-C222-pyant}



\begin{quote}
Before the new one: say the Kannada for a watermelon, then the Kannada for a chapati.
\end{quote}

\subsection*{You'll want to know: \kn{ಪ್ಯಾಂಟ್}}
\textbf{\kn{ಪ್ಯಾಂಟ್}} — \emph{pyāṇṭ} — \textquotedblleft{}trousers\textquotedblright{}.

\begin{grammarlens}[title={What you've built}]
One more word for this chapter.
\end{grammarlens}

\subsection*{Your turn}
\begin{itemize}
\raggedright
  \item \emph{Say it:} \emph{pyāṇṭ}
  \item \emph{Say it:} \emph{pyāṇṭ}, once more
  \item \emph{Say it:} say \emph{pyāṇṭ}
  \item \emph{Recall:} say the Kannada for a watermelon, then the Kannada for a chapati, then say \emph{pyāṇṭ} again
\end{itemize}

\begin{tcolorbox}[breakable,colback=teal!4,colframe=teal!35!black,title={Before you move on}]
What does \kn{ಪ್ಯಾಂಟ್} mean? (\textquotedblleft{}Trousers\textquotedblright{}.) Say it once more.
\end{tcolorbox}
\section[\texorpdfstring{\kn{ಲಂಗ} (laṅga)}{ಲಂಗ (laṅga)}]{\kn{ಲಂಗ} (laṅga) — a long skirt}

\label{lesson:KA-C222-langa}



\begin{quote}
Before the new one: say the Kannada for a chapati, then the Kannada for trousers.
\end{quote}

\subsection*{You'll want to know: \kn{ಲಂಗ}}
\textbf{\kn{ಲಂಗ}} — \emph{laṅga} — \textquotedblleft{}a long skirt\textquotedblright{}.

\begin{grammarlens}[title={What you've built}]
One more word for this chapter.
\end{grammarlens}

\subsection*{Your turn}
\begin{itemize}
\raggedright
  \item \emph{Say it:} \emph{laṅga}
  \item \emph{Say it:} \emph{laṅga}, once more
  \item \emph{Say it:} say \emph{laṅga}
  \item \emph{Recall:} say the Kannada for a chapati, then the Kannada for trousers, then say \emph{laṅga} again
\end{itemize}

\begin{tcolorbox}[breakable,colback=teal!4,colframe=teal!35!black,title={Before you move on}]
What does \kn{ಲಂಗ} mean? (\textquotedblleft{}A long skirt\textquotedblright{}.) Say it once more.
\end{tcolorbox}
\section[\texorpdfstring{\kn{ರವಿಕೆ} (ravike)}{ರವಿಕೆ (ravike)}]{\kn{ರವಿಕೆ} (ravike) — a blouse}

\label{lesson:KA-C222-ravike}



\begin{quote}
Before the new one: say the Kannada for trousers, then the Kannada for a long skirt.
\end{quote}

\subsection*{You'll want to know: \kn{ರವಿಕೆ}}
\textbf{\kn{ರವಿಕೆ}} — \emph{ravike} — \textquotedblleft{}a blouse\textquotedblright{}.

\begin{grammarlens}[title={What you've built}]
One more word for this chapter.
\end{grammarlens}

\subsection*{Your turn}
\begin{itemize}
\raggedright
  \item \emph{Say it:} \emph{ravike}
  \item \emph{Say it:} \emph{ravike}, once more
  \item \emph{Say it:} say \emph{ravike}
  \item \emph{Recall:} say the Kannada for trousers, then the Kannada for a long skirt, then say \emph{ravike} again
\end{itemize}

\begin{tcolorbox}[breakable,colback=teal!4,colframe=teal!35!black,title={Before you move on}]
What does \kn{ರವಿಕೆ} mean? (\textquotedblleft{}A blouse\textquotedblright{}.) Say it once more.
\end{tcolorbox}
\section[\texorpdfstring{\kn{ಬನಿಯನ್} (baniyan)}{ಬನಿಯನ್ (baniyan)}]{\kn{ಬನಿಯನ್} (baniyan) — a vest, an undershirt}

\label{lesson:KA-C222-baniyan}



\begin{quote}
Before the new one: say the Kannada for a long skirt, then the Kannada for a blouse.
\end{quote}

\subsection*{You'll want to know: \kn{ಬನಿಯನ್}}
\textbf{\kn{ಬನಿಯನ್}} — \emph{baniyan} — \textquotedblleft{}a vest, an undershirt\textquotedblright{}.

\begin{grammarlens}[title={What you've built}]
One more word for this chapter.
\end{grammarlens}

\subsection*{Your turn}
\begin{itemize}
\raggedright
  \item \emph{Say it:} \emph{baniyan}
  \item \emph{Say it:} \emph{baniyan}, once more
  \item \emph{Say it:} say \emph{baniyan}
  \item \emph{Recall:} say the Kannada for a long skirt, then the Kannada for a blouse, then say \emph{baniyan} again
\end{itemize}

\begin{tcolorbox}[breakable,colback=teal!4,colframe=teal!35!black,title={Before you move on}]
What does \kn{ಬನಿಯನ್} mean? (\textquotedblleft{}A vest, an undershirt\textquotedblright{}.) Say it once more.
\end{tcolorbox}
\section[\texorpdfstring{\kn{ಬೂಟು} (būṭu)}{ಬೂಟು (būṭu)}]{\kn{ಬೂಟು} (būṭu) — shoes}

\label{lesson:KA-C222-butu}



\begin{quote}
Before the new one: say the Kannada for a blouse, then the Kannada for a vest.
\end{quote}

\subsection*{You'll want to know: \kn{ಬೂಟು}}
\textbf{\kn{ಬೂಟು}} — \emph{būṭu} — \textquotedblleft{}shoes\textquotedblright{}.

\begin{grammarlens}[title={What you've built}]
One more word for this chapter.
\end{grammarlens}

\subsection*{Your turn}
\begin{itemize}
\raggedright
  \item \emph{Say it:} \emph{būṭu}
  \item \emph{Say it:} \emph{būṭu}, once more
  \item \emph{Say it:} say \emph{būṭu}
  \item \emph{Recall:} say the Kannada for a blouse, then the Kannada for a vest, then say \emph{būṭu} again
\end{itemize}

\begin{tcolorbox}[breakable,colback=teal!4,colframe=teal!35!black,title={Before you move on}]
What does \kn{ಬೂಟು} mean? (\textquotedblleft{}Shoes\textquotedblright{}.) Say it once more.
\end{tcolorbox}
